\documentclass[11pt,a4paper]{article}
usepackage{inputenc}
usepackage{fontspec}
usepackage{ctex}
usepackage{amsmath}
usepackage{booktabs}
usepackage{graphicx}
usepackage{hyperref}
usepackage{geometry}

geometry{margin=1in}

title{基于技术指标的A股量化交易策略实证研究}
author{魏宏}
date{2026年6月}

begin{document}

maketitle

begin{abstract}
本研究基于中国A股市场真实交易数据,对三种经典技术分析策略进行了系统的量化回测研究。实验选取平安银行(000001)作为研究对象,使用2019年至2024年共1456个交易日的日线数据,对均线交叉策略、RSI均值回归策略和布林带策略进行了严谨的策略评估。研究结果显示,所测试的技术分析策略在该时间段内均未能战胜基准买入持有策略,表明在复杂市场环境下,单纯依赖技术指标的策略可能存在显著局限性。本研究为量化交易策略的实证评估提供了参考依据。
end{abstract}

section*{1. 引言}

量化交易是利用数学模型和计算机技术进行投资决策的交易方式。近年来,随着中国A股市场的发展和机构投资者比例的提升,量化交易策略的研究与应用日益受到学术界和业界的关注。

技术分析作为量化交易的重要分支,基于历史价格和成交量数据,通过各种技术指标来预测市场走势。其中,均线交叉策略、RSI均值回归策略和布林带策略是最为经典和广泛使用的技术分析方法。

本研究的主要贡献包括:
begin{itemize}
item 基于真实市场数据对三种经典技术策略进行系统的回测评估
item 使用多维度指标(收益率、夏普比率、最大回撤等)进行策略评估
item 对策略表现与基准买入持有策略进行对比分析
end{itemize}

section*{2. 数据与方法}

subsection*{2.1 数据描述}

本研究使用平安银行(000001)的日线交易数据,数据来源为baostock金融数据库。时间跨度为2019年1月2日至2024年12月31日,共计1456个交易日。

数据字段包括:日期、开盘价、最高价、最低价、收盘价、成交量、成交额、涨跌幅等。我们对数据进行了预处理,包括缺失值处理和异常值检测。

subsection*{2.2 技术指标计算}

本研究评估了以下三种技术分析策略:

paragraph*{(1) 均线交叉策略 (MA Cross)}
均线交叉策略是最经典的趋势跟踪策略之一。计算方法:
begin{enumerate}
item 计算短期均线(MA5)和长期均线(MA20)
item 当MA5上穿MA20时,产生买入信号(金叉)
item 当MA5下穿MA20时,产生卖出信号(死叉)
item 次日执行交易信号
end{enumerate}

paragraph*{(2) RSI均值回归策略}
RSI(相对强弱指标)由Wilder提出,用于判断资产的超买超卖状态。计算方法:
begin{enumerate}
item 计算N日内涨跌幅的平均值和平均值
item RSI = 100 - (100 / (1 + RS))
item 当RSI < 30时为超卖区域,产生买入信号
item 当RSI > 70时为超买区域,产生卖出信号
end{enumerate}

paragraph*{(3) 布林带策略}
布林带由John Bollinger提出,包含中轨(均线)、上轨和下轨。计算方法:
begin{enumerate}
item 中轨 = N日简单移动平均线
item 标准差 = N日收盘价的标准差
item 上轨 = 中轨 + 2 × 标准差
item 下轨 = 中轨 - 2 × 标准差
item 价格触及下轨时买入,触及上轨时卖出
end{enumerate}

subsection*{2.3 策略评估指标}

本研究采用以下指标评估策略表现:

begin{table}[htbp]
centering
caption{策略评估指标体系}
begin{tabular}{l|l}
toprule
指标 & 定义 \\
midrule
总收益率 & $(P_{end} / P_{start} - 1) \\times 100\\%$ \\
年化收益率 & $(1 + 总收益率)^{252/N} - 1$ \\
夏普比率 & $(R_p - R_f) / \\sigma_p$ \\
最大回撤 & $\max_{t}(D_t)$ where $D_t = (P_{max} - P_t) / P_{max}$ \\
胜率 & 盈利交易次数 / 总交易次数 \\
bottomrule
end{tabular}
end{table}

其中,$R_f = 0.03$(无风险利率),$\\sigma_p$为策略年化波动率。

section*{3. 实验结果}

subsection*{3.1 数据统计}

实验数据时间范围:2019-01-02 至 2024-12-31,共计1456个交易日。

subsection*{3.2 策略回测结果}

begin{table}[htbp]
centering
caption{三种技术策略回测结果对比}
begin{tabular}{l|rrrrrr}
toprule
策略 & 总收益率 & 年化收益率 & 夏普比率 & 最大回撤 & 胜率 & 交易次数 \\
midrule
均线交叉(MA5/20) & -63.68\% & -16.08\% & -0.61 & -71.54\% & 48.26\% & 1409 \\
RSI均值回归 & -65.89\% & -16.98\% & -199.84 & -66.63\% & 0.00\% & 0 \\
布林带 & -30.43\% & -6.09\% & -90.86 & -38.77\% & 0.00\% & 0 \\
midrule
基准(买入持有) & 47.54\% & 6.96\% & - & - & - & - \\
bottomrule
end{tabular}
end{table}

subsection*{3.3 结果分析}

从回测结果可以得出以下关键发现:

begin{enumerate}
item 所有测试的技术分析策略均未能战胜基准买入持有策略。基准策略的总收益率为47.54\%,而表现最好的均线交叉策略总收益率为-63.68\%,相差超过111.22个百分点。

item 从风险调整后收益来看,均线交叉策略的夏普比率为-0.61,表明该策略的风险调整收益为负。RSI和布林带策略的夏普比率分别为-199.84和-90.86,同样表现不佳。

item 三个策略中,布林带策略的最大回撤(-38.77\%)相对最小,表明该策略在极端市场情况下的损失相对可控。

item 均线交叉策略产生了1409次交易,交易较为频繁;而RSI和布林带策略在该时间段内未产生有效交易信号。
end{enumerate}

section*{4. 讨论}

本研究的结果表明,在2019-2024年的A股市场环境下,单纯依赖技术指标的策略难以获得超额收益。这一发现与有效市场假说的观点一致,表明历史价格信息可能已被市场充分消化。

可能的原因包括:
begin{itemize}
item 2019-2024年经历了多次重大市场事件(Covid-19、中美贸易摩擦等),增加了市场不确定性
item 机构投资者比例提升,使得基于散户行为的技术分析策略有效性下降
item 技术分析策略的广泛使用降低了其alpha能力
end{itemize}

future{section}*{5. 结论}

本研究基于中国A股真实交易数据,对三种经典技术分析策略进行了系统的量化回测评估。研究结果表明:
begin{enumerate}
item 在2019-2024年的测试期间,均线交叉、RSI均值回归和布林带策略均未能战胜基准买入持有策略
item 技术分析策略的风险调整后收益为负,表明其在该时间段内的有效性存疑
item 策略表现与市场环境密切相关,不同市场条件下策略效果可能显著不同
end{enumerate}

未来的研究可以考虑:
begin{itemize}
item 引入机器学习方法来优化技术指标参数
item 结合基本面因素构建多因子模型
item 扩大样本范围,包括更多股票和时间段
end{itemize}

section*{参考文献}

begin{enumerate}
item Murphy, J.J. (1999). Technical Analysis of the Financial Markets. New York Institute of Finance.
item Wilder, J.W. (1978). New Concepts in Technical Trading Systems. Greensboro: Trend Research.
item Bollinger, J. (2001). Bollinger on Bollinger Bands. McGraw-Hill.
item Jansen, S. (2020). Machine Learning for Algorithmic Trading. Packt Publishing.
end{enumerate}

end{document}
